\documentclass[sigplan,screen,anonymous,nonacm]{acmart}

\usepackage{mathpartir}
\usepackage{booktabs}
\usepackage[nameinlink,noabbrev]{cleveref}

\settopmatter{printacmref=false,printccs=false,printfolios=true}

\newcommand{\lambdap}{\ensuremath{\lambda_{\mathsf p}}}
\newcommand{\ctxempty}{\varnothing}
\newcommand{\sing}[1]{\{#1\}}
\newcommand{\pairty}[4]{\langle #1\!:\!#2;\,#3\!:\!#4\rangle}
\newcommand{\pairval}[3]{\langle #1;\,#2=#3\rangle}
\newcommand{\subst}[3]{#1[#2/#3]}
\newcommand{\wf}{\mathrel{\mathsf{wf}}}
\newcommand{\initial}{\mathsf{initial}}
\newcommand{\final}{\mathsf{final}}
\newcommand{\steps}{\longrightarrow^{*}}
\newcommand{\sstep}{\longrightarrow}
\newcommand{\resolve}{\Downarrow}
\newcommand{\semmap}{\mathrel{\Rightarrow_{\!s}}}
\newcommand{\runtimeeq}{\mathrel{\equiv_{\sigma}}}
\newcommand{\struct}{\Vdash}
\newcommand{\exact}{\mathsf{exact}}

\title{\texorpdfstring{\lambdap}{lambda-p}: Paths through Dependent Pairs}
\author{Anonymous Author(s)}
\renewcommand{\shortauthors}{Anonymous Author(s)}

\begin{document}

\begin{abstract}
We present a monadic-normal-form calculus in which paths traverse immutable
chains of dependent pairs. Pair definitions store either a term reference or
a type, and the member signature may refer to the pair's first component.
The static semantics distinguishes term singletons from selected abstract
type members. If singleton symmetry is instead allowed to treat abstract
membership as path identity, a closed, well-typed program becomes stuck. The
calculus uses a conservative subtyping rule for dependent pairs.
It permits member types to change when the first component has singleton type,
and permits the first component to be widened separately. We prove type safety
in Lean. The proof records exact value-introduction types in the store and
interprets structural subtyping by finite, proof-relevant coercion codes.
Every endpoint of a finite execution of a closed, well-typed term is final or
can take another step.
\end{abstract}

\keywords{path-dependent types, dependent pairs, singleton types, type
soundness, Lean}

\maketitle

\section{Introduction}

A path-dependent type assigns a type component to a stable term path. The
usual object-calculus presentation combines this mechanism with recursive
objects, refinements, intersections, and a substantial subtyping relation
\citep{amin2014foundations,rompf2016dot}. The purpose of \lambdap{} is
narrower: to isolate paths, abstract type members, and dependency between the
two components of a pair.

Values in \lambdap{} are functions and immutable dependent pairs. A pair
contains a first component and either a term reference or a type definition. Its
member signature may mention the first component. Paths follow first
projections and named members through chains of pairs; a missed label resumes
the search in the first component. The static and dynamic rules expose this
path mechanism directly.

Two design points are essential to the result in this paper. First, the
static semantics distinguishes the singleton type \(\sing{p}\), which states
that a term denotes the same value as \(p\), from the abstract type selection
\(p.A\). This calculus makes the distinction syntactic. An earlier version
used one constructor for both roles. We give a closed counterexample to that
version: singleton symmetry turns membership in an abstract type into a false
path equality.

Second, dependent-pair subtyping is deliberately conservative. One rule
widens the first component without changing the member. A second rule changes
the member only when the first component is the singleton of a path, and
retains an explicit subtyping premise after opening the member at that path.
The rules suffice to hide exact type definitions behind dependent abstract
bounds. We do not know whether the natural unrestricted pair rule is unsound;
the restriction records the evidence required by the direct semantic proof.

The main technical result is a Lean-checked finite-run safety theorem. The
proof has three parts.

\begin{itemize}
\item An exact store invariant records every allocated value at its
  syntax-directed introduction type.
\item A store-indexed path equivalence accounts for paths that resolve to the
  same location. A proof-only structural typing layer admits conversion by
  this equivalence.
\item A finite syntax of semantic coercions interprets every structural
  subtyping derivation. Applying a coercion to a location singleton yields
  the canonical forms used by progress and preservation.
\end{itemize}

The result concerns a deliberately limited language. There are no recursive
objects, recursive allocation, mutation, intersections, unions, or general
source terms outside monadic normal form. We prove neither normalization nor
decidability, and declarative transitivity remains a primitive rule.

\section{The calculus}
\label{sec:calculus}

\subsection{Syntax and binding}

\Cref{fig:syntax} gives the syntax. Kinds distinguish term members
\((\star)\) from type members \((\iota)\). A generalized type
\(\tau^\kappa\) is a proper type at kind \(\star\), and an interval at kind
\(\iota\). A definition has the corresponding kind: a term member stores a
variable, while a type member stores a proper type.

\begin{figure}[t]
\scriptsize
\[
\begin{array}{rcl}
\kappa &::=& \star \mid \iota\\
p,q &::=& x \mid p.1 \mid p.a\\
T,U &::=& \top \mid \bot \mid (x:S)\to T
  \mid \pairty{x}{S}{a}{\tau^\kappa}\\
&& {}\mid \sing{p} \mid p.A\\
\tau^\star &::=& T\\
\tau^\iota &::=& L..U\\
\delta^\star &::=& z\\
\delta^\iota &::=& T\\
t,u &::=& p \mid \lambda(x:S).t
  \mid \pairval{y}{a}{\delta^\kappa}\\
&& {}\mid p\,q
  \mid \mathbf{let}\ x=t\ \mathbf{in}\ u
  \mid (t:T)\\
v &::=& \lambda(x:S).t \mid \pairval{y}{a}{\delta^\kappa}\\
\Gamma &::=& \ctxempty \mid \Gamma,x:T
\end{array}
\]
\caption{Source syntax. Uppercase labels conventionally denote type members.}
\Description{Grammar of kinds, paths, proper and generalized types,
definitions, terms, values, and contexts.}
\label{fig:syntax}
\end{figure}

The binder \(x\) in a function scopes over its codomain. In
\(\pairty{x}{S}{a}{\tau}\), it scopes over the member signature \(\tau\).
We write \(\subst{\tau}{p}{x}\) for capture-avoiding opening. The
mechanization uses intrinsically scoped de Bruijn syntax; the paper uses
names.

Terms are in monadic normal form. Pair components are variables, and the
operator and operand of an application are paths. This restriction leaves
only \(\mathbf{let}\) frames in the evaluation context and makes allocation
explicit.

The four source judgments are
\[
\Gamma\vdash p::\tau^\kappa
\qquad
\Gamma\vdash\tau_1^\kappa<:\tau_2^\kappa
\qquad
\Gamma\vdash\tau^\kappa\ \wf
\qquad
\Gamma\vdash t:T.
\]
The double colon denotes \emph{precise path typing}: unlike term typing, it
does not contain subsumption.

\subsection{Paths and pair chains}

The precise path rules appear in the upper half of
\cref{fig:path-wf}. A variable synthesizes its context entry. A first
projection synthesizes the first-component type. Selecting the matching
member opens its signature with the first projection of the receiver.

Pairs also represent record chains. If the outer label is \(b\) but the
requested label is \(a\), rule \textsc{P-Miss} continues from \(p.1.a\).
This static rule matches the recursive missed-label case of path resolution
in \cref{fig:dynamics}.

\begin{figure*}[t]
\begin{minipage}{\textwidth}
\footnotesize
\begin{mathpar}
\inferrule*[right=\textsc{P-Var}]
  {\Gamma(x)=T}
  {\Gamma\vdash x::T}
\and
\inferrule*[right=\textsc{P-Fst}]
  {\Gamma\vdash p::\pairty{x}{S}{a}{\tau}}
  {\Gamma\vdash p.1::S}
\and
\inferrule*[right=\textsc{P-Sel}]
  {\Gamma\vdash p::\pairty{x}{S}{a}{\tau}}
  {\Gamma\vdash p.a::\subst{\tau}{p.1}{x}}
\and
\inferrule*[right=\textsc{P-Miss}]
  {\Gamma\vdash p::\pairty{x}{S}{b}{\tau'}
   \\
   \Gamma\vdash p.1.a::\tau
   \\
   a\ne b}
  {\Gamma\vdash p.a::\tau}
\end{mathpar}

\vspace{-0.5\baselineskip}
\hrulefill
\vspace{-0.5\baselineskip}

\begin{mathpar}
\inferrule*[right=\textsc{W-Bot}]
  { }{\Gamma\vdash\bot\ \wf}
\and
\inferrule*[right=\textsc{W-Top}]
  { }{\Gamma\vdash\top\ \wf}
\and
\inferrule*[right=\textsc{W-Single}]
  {\Gamma\vdash p::T}
  {\Gamma\vdash\sing{p}\ \wf}
\and
\inferrule*[right=\textsc{W-Select}]
  {\Gamma\vdash p::\pairty{x}{S}{A}{L..U}}
  {\Gamma\vdash p.A\ \wf}
\and
\inferrule*[right=\textsc{W-Fun}]
  {\Gamma\vdash S\ \wf
   \\
   \Gamma,x:S\vdash T\ \wf}
  {\Gamma\vdash(x:S)\to T\ \wf}
\and
\inferrule*[right=\textsc{W-Pair}]
  {\Gamma\vdash S\ \wf
   \\
   \Gamma,x:S\vdash\tau\ \wf}
  {\Gamma\vdash\pairty{x}{S}{a}{\tau}\ \wf}
\and
\inferrule*[right=\textsc{W-Bounds}]
  {\Gamma\vdash L\ \wf
   \\
   \Gamma\vdash U\ \wf
   \\
   \Gamma\vdash L<:U}
  {\Gamma\vdash L..U\ \wf}
\end{mathpar}
\end{minipage}
\caption{Precise path typing (above) and well-formedness (below).}
\Description{Inference rules for precise path typing and generalized type
well-formedness.}
\label{fig:path-wf}
\end{figure*}

\subsection{Subtyping and well-formedness}

\Cref{fig:subtyping} gives the complete declarative subtyping relation.
It includes reflexivity and primitive transitivity, as well as the usual
extremal and function rules. Rule \textsc{S-Widen} exposes the precise type
of a singleton path. Rule \textsc{S-Alias} is singleton symmetry: if \(p\)
has the \emph{term singleton} \(\sing{q}\), then the singleton subtyping
relation may be reversed.

Abstract type selections use the interval synthesized by precise member
lookup. The premise \(L<:U\) is an explicit nonemptiness guard. It is retained
both here and in interval formation and subtyping. Interval subtyping is
contravariant in its lower bound and covariant in its upper bound.

\begin{figure*}[t]
\begin{minipage}{\textwidth}
\footnotesize
\begin{mathpar}
\inferrule*[right=\textsc{S-Refl}]
  { }{\Gamma\vdash\tau<:\tau}
\and
\inferrule*[right=\textsc{S-Trans}]
  {\Gamma\vdash\tau_1<:\tau_2
   \\
   \Gamma\vdash\tau_2<:\tau_3}
  {\Gamma\vdash\tau_1<:\tau_3}
\and
\inferrule*[right=\textsc{S-Bot}]
  { }{\Gamma\vdash\bot<:T}
\and
\inferrule*[right=\textsc{S-Top}]
  { }{\Gamma\vdash T<:\top}
\and
\inferrule*[right=\textsc{S-Widen}]
  {\Gamma\vdash p::T}
  {\Gamma\vdash\sing{p}<:T}
\and
\inferrule*[right=\textsc{S-Alias}]
  {\Gamma\vdash p::\sing{q}}
  {\Gamma\vdash\sing{q}<:\sing{p}}
\and
\inferrule*[right=\textsc{S-Sel-Lo}]
  {\Gamma\vdash p.A::L..U
   \\
   \Gamma\vdash L<:U}
  {\Gamma\vdash L<:p.A}
\and
\inferrule*[right=\textsc{S-Sel-Hi}]
  {\Gamma\vdash p.A::L..U
   \\
   \Gamma\vdash L<:U}
  {\Gamma\vdash p.A<:U}
\and
\inferrule*[right=\textsc{S-Fun}]
  {\Gamma\vdash S'<:S
   \\
   \Gamma,x:S'\vdash T<:T'}
  {\Gamma\vdash (x:S)\to T<:(x:S')\to T'}
\and
\inferrule*[right=\textsc{S-Pair-Fst}]
  {\Gamma\vdash S<:S'}
  {\Gamma\vdash
    \pairty{x}{S}{a}{\tau}
    <:
    \pairty{x}{S'}{a}{\tau}}
\and
\inferrule*[right=\textsc{S-Pair-Single}]
  {\Gamma\vdash p::P
   \\
   \Gamma,x:\sing{p}\vdash\tau<:\tau'
   \\
   \Gamma\vdash\subst{\tau}{p}{x}<:\subst{\tau'}{p}{x}}
  {\Gamma\vdash
    \pairty{x}{\sing{p}}{a}{\tau}
    <:
    \pairty{x}{\sing{p}}{a}{\tau'}}
\and
\inferrule*[right=\textsc{S-Bounds}]
  {\Gamma\vdash L'<:L
   \\
   \Gamma\vdash U<:U'
   \\
   \Gamma\vdash L<:U}
  {\Gamma\vdash L..U<:L'..U'}
\end{mathpar}
\end{minipage}
\caption{Generalized subtyping. Premises and conclusions have the same kind.}
\Description{The complete declarative generalized subtyping relation,
including primitive transitivity, singleton and selection rules, functions,
the two restricted pair rules, and interval subtyping.}
\label{fig:subtyping}
\end{figure*}

The two pair rules separate two operations. \textsc{S-Pair-Fst} changes the
first-component type and leaves the member fixed. \textsc{S-Pair-Single}
changes the member under a singleton first component. Its last premise is not a
well-formedness condition: it is a second subtyping derivation after opening
the member at \(p\). Declarative transitivity composes the two operations.
\Cref{sec:pairs} explains the restriction and its use in the proof.

The lower half of \cref{fig:path-wf} gives all well-formedness rules.
Singleton formation requires a term-denoting path, while abstract selection
requires a pair with an interval member. Functions and pairs check their
dependent component under the corresponding binder. Rule
\textsc{W-Bounds} repeats the nonemptiness condition.

\subsection{Term typing}

\Cref{fig:typing} gives term typing. A path receives its singleton type;
\textsc{S-Widen} can subsequently expose its precise type. Application is
dependent. The two pair-introduction rules assign exact member signatures.
A term definition receives \(\sing{z}\); a type definition receives \(T..T\).
The first component is also exact, namely \(\sing{y}\).

The result type of a let expression cannot depend on its bound variable. The
side condition is represented in the intrinsically scoped development by
weakening \(T\) below the binder. Subsumption requires the target to be
well-formed.

\begin{figure*}[t]
\begin{minipage}{\textwidth}
\footnotesize
\begin{mathpar}
\inferrule*[right=\textsc{T-Path}]
  {\Gamma\vdash p::T}
  {\Gamma\vdash p:\sing{p}}
\and
\inferrule*[right=\textsc{T-Abs}]
  {\Gamma,x:S\vdash t:T
   \\
   \Gamma\vdash S\ \wf}
  {\Gamma\vdash\lambda(x:S).t:(x:S)\to T}
\and
\inferrule*[right=\textsc{T-App}]
  {\Gamma\vdash p:(x:S)\to T
   \\
   \Gamma\vdash q:S}
  {\Gamma\vdash p\,q:\subst{T}{q}{x}}
\and
\inferrule*[right=\textsc{T-Pair}]
  {\Gamma(y)=S
   \\
   \Gamma(z)=T}
  {\Gamma\vdash
    \pairval{y}{a}{z}:
    \pairty{x}{\sing{y}}{a}{\sing{z}}}
\and
\inferrule*[right=\textsc{T-Type-Pair}]
  {\Gamma(y)=S
   \\
   \Gamma\vdash T\ \wf}
  {\Gamma\vdash
    \pairval{y}{A}{T}:
    \pairty{x}{\sing{y}}{A}{T..T}}
\and
\inferrule*[right=\textsc{T-Let}]
  {\Gamma\vdash s:S
   \\
   \Gamma\vdash T\ \wf
   \\
   \Gamma,x:S\vdash t:T
   \\
   x\notin\mathsf{fv}(T)}
  {\Gamma\vdash
    \mathbf{let}\ x=s\ \mathbf{in}\ t:T}
\and
\inferrule*[right=\textsc{T-Ascribe}]
  {\Gamma\vdash t:T
   \\
   \Gamma\vdash T\ \wf}
  {\Gamma\vdash(t:T):T}
\and
\inferrule*[right=\textsc{T-Sub}]
  {\Gamma\vdash t:S
   \\
   \Gamma\vdash S<:T
   \\
   \Gamma\vdash T\ \wf}
  {\Gamma\vdash t:T}
\end{mathpar}
\end{minipage}
\caption{Term typing. Names bound by pair types are in scope in their member
signature.}
\Description{Typing rules for paths, abstractions, applications, term and
type pairs, let expressions, ascriptions, and subsumption.}
\label{fig:typing}
\end{figure*}

\subsection{Store and machine}

A store \(\sigma\) is an immutable sequence of values. A stored value may
refer only to earlier cells. We write \(\sigma(x)=v\) for lookup and
\(\sigma,x\mapsto v\) for extension at a fresh location. There is one
intrinsically scoped namespace for source variables and run-time locations;
the named presentation calls store-bound variables locations when useful.

Source path resolution, written \(\sigma\vdash p\resolve x\), is big-step.
The rules are at the top of \cref{fig:dynamics}. A variable resolves to
itself. Projection returns the stored first component. A matching term
member returns its stored location. A missed label continues at the first
component. Resolution is deterministic, although it is partial on arbitrary
stores and paths.

Machine states are triples \(\langle\sigma,K,t\rangle\). Continuations
contain only let frames \(F=\mathbf{let}\ x=[\,]\ \mathbf{in}\ t\). The six
transition rules appear at the bottom of \cref{fig:dynamics}. Application
resolves both paths before opening a stored abstraction. A non-variable path
normalizes to its resolved location. A returned location is substituted
directly into a suspended let body; a syntactic value is instead allocated
at a fresh location. Store extension weakens the old continuation, written
\(K^\uparrow\), into the extended scope. Ascriptions have no dynamic effect.

\begin{figure*}[t]
\begin{minipage}{\textwidth}
\footnotesize
\begin{mathpar}
\inferrule*[right=\textsc{R-Var}]
  { }{\sigma\vdash x\resolve x}
\and
\inferrule*[right=\textsc{R-Fst}]
  {\sigma\vdash p\resolve x
   \\
   \sigma(x)=\pairval{y}{a}{\delta}}
  {\sigma\vdash p.1\resolve y}
\and
\inferrule*[right=\textsc{R-Hit}]
  {\sigma\vdash p\resolve x
   \\
   \sigma(x)=\pairval{y}{a}{z}}
  {\sigma\vdash p.a\resolve z}
\and
\inferrule*[right=\textsc{R-Miss}]
  {\sigma\vdash p\resolve x
   \\
   \sigma(x)=\pairval{y}{b}{\delta}
   \\
   a\ne b
   \\
   \sigma\vdash y.a\resolve z}
  {\sigma\vdash p.a\resolve z}
\end{mathpar}

\vspace{-0.5\baselineskip}
\hrulefill
\vspace{-0.5\baselineskip}

\begin{mathpar}
\inferrule*[right=\textsc{E-App}]
  {\sigma\vdash p\resolve f
   \\
   \sigma\vdash q\resolve y
   \\
   \sigma(f)=\lambda(x:S).t}
  {\langle\sigma,K,p\,q\rangle
   \sstep
   \langle\sigma,K,\subst{t}{y}{x}\rangle}
\and
\inferrule*[right=\textsc{E-Path}]
  {\sigma\vdash p\resolve x
   \\
   p\ \mathsf{not\ a\ variable}}
  {\langle\sigma,K,p\rangle
   \sstep
   \langle\sigma,K,x\rangle}
\and
\inferrule*[right=\textsc{E-Let}]
  {F=\mathbf{let}\ x=[\,]\ \mathbf{in}\ t}
  {\langle\sigma,K,\mathbf{let}\ x=s\ \mathbf{in}\ t\rangle
   \sstep
   \langle\sigma,F::K,s\rangle}
\and
\inferrule*[right=\textsc{E-Return}]
  {F=\mathbf{let}\ x=[\,]\ \mathbf{in}\ t}
  {\langle\sigma,F::K,y\rangle
   \sstep
   \langle\sigma,K,\subst{t}{y}{x}\rangle}
\and
\inferrule*[right=\textsc{E-Alloc}]
  {F=\mathbf{let}\ x=[\,]\ \mathbf{in}\ t
   \\
   v\ \mathsf{value}
   \\
   \ell\ \mathsf{fresh}}
  {\langle\sigma,F::K,v\rangle
   \sstep
   \langle\sigma[\ell\mapsto v],K^\uparrow,
     \subst{t}{\ell}{x}\rangle}
\and
\inferrule*[right=\textsc{E-Ascribe}]
  { }
  {\langle\sigma,K,(t:T)\rangle
   \sstep
   \langle\sigma,K,t\rangle}
\end{mathpar}
\end{minipage}
\caption{Big-step path resolution (above) and CK-machine transitions
(below).}
\Description{Rules for resolving paths through an immutable store and for
stepping applications, paths, let frames, allocations, and ascriptions.}
\label{fig:dynamics}
\end{figure*}

A state is final when its continuation is empty and its term is either a
valid store location or a syntactic value. The initial state of \(t\) is
\(\initial(t)=\langle\varnothing,[\,],t\rangle\).

\section{Singletons, selections, and pair abstraction}
\label{sec:pairs}

\subsection{Why singleton types and selections are distinct}

The term singleton \(\sing{p}\) asserts value identity. In particular,
\textsc{S-Alias} is justified when precise typing says that one path has
another path's singleton type. An abstract selection \(p.A\), by contrast,
denotes an unknown type between the bounds stored at \(A\). Membership in
\(p.A\) says nothing by itself about path identity.

Some distinction between the two roles is required by a source-level
counterexample. In a precursor to \lambdap{}, both forms were represented by
a single type constructor and singleton symmetry accepted either role. Let
\[
\begin{array}{rcl}
f &=& \lambda(x:\top).x,\\
q &=& \pairval{f}{A}{\top},\\
h &=& \lambda(z:q.A).q.
\end{array}
\]
The historical rules type the body of \(h\) by the chain
\[
  \sing{q}<:\top<:q.A<:\sing{z}.
\]
The last edge applies singleton symmetry to the context binding \(z:q.A\),
treating abstract membership as equality. The same rules accept \(f\) as an
argument of type \(q.A\). Thus \(h\,f\) is assigned \(\sing{f}\), although it
returns \(q\).

The closed program
\[
\begin{array}{l}
\mathbf{let}\ f=\lambda(x:\top).x\ \mathbf{in}\\
\mathbf{let}\ q=\pairval{f}{A}{\top}\ \mathbf{in}\\
\mathbf{let}\ h=\lambda(z:q.A).q\ \mathbf{in}\\
\mathbf{let}\ r=h\,f\ \mathbf{in}\ r\,f
\end{array}
\]
takes nine machine steps to an application whose operator location contains
the pair \(q\). It is neither final nor reducible. Every interval in the
derivation is \(\top..\top\), so this is not a bad-bounds example. It also
does not depend on transitivity being primitive: any ordinary transitivity
lemma composes the displayed chain.

The full derivation and trace are retained in the historical library, in
module \texttt{SourceUnsoundnessCounterexample}.
In the present syntax, a binding \(z:q.A\) cannot be a premise for
\textsc{S-Alias}. Theorems
\nolinkurl{historical_body_subtyping_blocked} and
\nolinkurl{selection_to_argument_singleton_blocked} show that the crucial
subtypings remain underivable even with arbitrary uses of primitive
transitivity.

The counterexample rules out treating arbitrary abstract membership as
singleton evidence. It does not show that two abstract-syntax constructors
are the only possible repair: one could instead add a term-denotation
classification premise to singleton symmetry. Separate constructors make the
distinction explicit and are the design adopted here.

\subsection{The pair rule}

The unrestricted rule one might first write is
\[
\inferrule
  {\Gamma\vdash S<:S'
   \\
   \Gamma,x:S\vdash\tau<:\tau'}
  {\Gamma\vdash
    \pairty{x}{S}{a}{\tau}
    <:
    \pairty{x}{S'}{a}{\tau'}}.
\tag{Unrestricted Pair}
\]
Its semantic interpretation must work for an arbitrary run-time first
component \(y\) that realizes \(S\). In particular, it must turn a realizer
of \(\subst{\tau}{y}{x}\) into a realizer of
\(\subst{\tau'}{y}{x}\). The raw subtyping derivation under \(x:S\) does not
contain a structurally smaller semantic coercion specialized to \(y\).
Representing the required inhabitant-indexed family directly as a function
would place the realization relation negatively in its own definition.

\textsc{S-Pair-Single} fixes a path \(p\) and supplies the specialized
subtyping derivation explicitly. If the stored first component \(y\)
realizes \(\sing{p}\), then \(y\) and \(p\) resolve to the same location.
Run-time conversion identifies
\[
\subst{\tau}{y}{x}
\runtimeeq
\subst{\tau}{p}{x},
\]
the explicit opened coercion acts at \(p\), and conversion transports the
result back to \(\subst{\tau'}{y}{x}\). This is the semantic step for which
the extra premise of \textsc{S-Pair-Single} is needed.

The restriction still supports the intended abstract interface. If
\(\Gamma\vdash p::P\), the checked examples derive
\[
\begin{aligned}
&\pairty{x}{\sing{p}}{A}{\sing{p}..\sing{p}}
\\[-1mm]
&\quad<:
\pairty{x}{\sing{p}}{A}{\bot..\sing{x}}
\\[-1mm]
&\quad<:
\pairty{x}{P}{A}{\bot..\sing{x}}.
\end{aligned}
\]
The first step changes an exact member into dependent abstract bounds while
the first component is a singleton. The second step widens the first
component. The order matters: after widening to an arbitrary \(P\), the
calculus does not generally permit changing the member.

This is a proof-motivated conservative restriction, not a demonstrated
boundary of soundness. We have no counterexample to the unrestricted rule
after singleton types and abstract selections have been separated.

\section{Soundness}
\label{sec:soundness}

\subsection{Statements}

Because allocation extends the intrinsically scoped store, finite machine
executions relate states at possibly different scope sizes. We write
\(\initial(t)\steps s\) for this heterogeneous reflexive-transitive closure.
Progress has its usual state formulation:
\[
\mathsf{Progress}(s)
\quad\stackrel{\mathrm{def}}{\Longleftrightarrow}\quad
\final(s)\ \lor\ \exists s'.\,s\sstep s'.
\]

\begin{theorem}[Finite-run type safety]
\label{thm:safety}
If \(\ctxempty\vdash t:T\) and \(\initial(t)\steps s\), then
\(\mathsf{Progress}(s)\).
\end{theorem}

The preservation invariant is more precise than source typing. Write
\(\Delta\struct_{\exact}s:U\) when the store is recorded at exact
value-introduction types and the current term and continuation are checked by
the structural run-time judgments described below. Write
\((\Gamma,T)\preceq(\Delta,U)\) when \(\Delta\) extends \(\Gamma\) by zero or
more allocations and \(U\) is weakened once for each extension.

\begin{theorem}[Finite-run preservation]
\label{thm:preservation}
If \(\ctxempty\vdash t:T\) and \(\initial(t)\steps s\), then there exist
\(\Delta\) and \(U\) such that
\[
  (\ctxempty,T)\preceq(\Delta,U)
  \qquad\text{and}\qquad
  \Delta\struct_{\exact}s:U.
\]
\end{theorem}

The source judgment is designed for source terms, not intermediate CK
states. \Cref{thm:preservation} states closure of the actual machine
invariant; \cref{thm:safety} is its source-facing non-stuckness corollary.

\subsection{Path resolution and exact stores}

The first difficulty is path lookup. Under a well-typed exact store, every
proper, precisely typed path resolves. Separately, resolution preserves
ordinary source typing of its endpoint. It does \emph{not} preserve the same
precise synthesized type literally: a projection can synthesize a singleton,
while the location it reaches synthesizes its context entry. The two useful
statements are therefore
\[
\begin{gathered}
\inferrule
  {\Gamma\vdash p::T
   \\
   \mathsf{ExactStore}(\Gamma,\sigma)}
  {\exists x.\ \sigma\vdash p\resolve x}
\\[1.2ex]
\inferrule
  {\Gamma\vdash p:T
   \\
   \mathsf{ExactStore}(\Gamma,\sigma)
   \\
   \sigma\vdash p\resolve x}
  {\Gamma\vdash x:T}
\end{gathered}
\]

An exact store assigns a cell the syntax-directed type of the value that was
allocated there. This differs from an ordinary store-typing relation that
allows a cell to be recorded at a supertype. Exactness is needed for
canonical forms: if a location is later viewed as a function, the proof must
recover the actual introduction form and its original signature. In the
allocation case, value inversion recovers the introduction type, and
narrowing rechecks the suspended let body under that exact type.

\subsection{Run-time conversion}

Different paths may reach the same location. Let \(p\runtimeeq q\) be the
least congruence containing every pair of paths that co-resolve in
\(\sigma\). Types and generalized types admit conversion along this
equivalence. Conversion is used, for example, to relate a selected member
opened at \(p.1\) to the same member opened at the location returned by the
projection.

Ordinary weakening is not sufficient below binders: the new variable has no
store cell and must not inherit accidental equations from the old scope.
The proof therefore uses a scoped lift of \(\runtimeeq\). Old equations are
weakened, while the new variable is related only to itself.

The source rules embed into proof-only \emph{structural} judgments for paths,
subtyping, well-formedness, terms, continuations, and states. These judgments
expose the constructors of the source derivation and add conversion by
\(\runtimeeq\). They do not change the source calculus; they provide the
induction principles needed for preservation.

\subsection{Generalized resolution and semantic maps}

Source evaluation resolves term paths only to locations. The proof uses a
generalized endpoint
\[
  e ::= \mathsf{val}(x) \mid \mathsf{type}(W),
\]
so one resolution relation covers both term selection and stored type
definitions. On term paths it agrees with source resolution.

The interpretation consists of three mutually inductive judgments:
proper-type possibility, generalized endpoint realization, and semantic
maps. A judgment \(\mathsf{Possible}_{\Gamma,\sigma}(x,T)\) says that
location \(x\) can be used at proper type \(T\). Its key cases are:

\begin{itemize}
\item a function type requires an actual stored abstraction and records the
  residual domain and codomain subtypings needed for beta reduction;
\item a pair type requires an actual stored pair and realizers for its first
  component and opened member;
\item \(\sing{p}\) requires \(p\) to resolve to \(x\); and
\item \(p.A\) requires \(p.A\) to resolve to a stored type \(W\), followed by
  a realizer of \(W\).
\end{itemize}

Endpoint realization handles generalized types. A value endpoint realizes a
proper type through \(\mathsf{Possible}\). A type endpoint \(W\) realizes
\(L..U\) by retaining a lower coercion \(L\semmap W\) and an upper coercion
\(W\semmap U\).

The notation \(\tau_1\semmap\tau_2\) denotes a finite, proof-relevant
\emph{semantic map}. Its constructors follow structural subtyping:
identity, composition, conversion, top and bottom, singleton widening and
aliasing, selection bounds, functions, the two pair rules, and intervals.
It is a syntax of coercions rather than a semantic function, so the mutual
definition remains strictly positive. Each code has an action on endpoint
realizers:
\[
\inferrule
  {\tau_1\semmap\tau_2
   \\
   e\in\llbracket\tau_1\rrbracket_{\Gamma,\sigma}}
  {e\in\llbracket\tau_2\rrbracket_{\Gamma,\sigma}}.
\tag{Map Action}
\]

\subsection{Fundamental lemma}

The central induction is stated under simultaneous path substitution. Call
\(\rho\) a \emph{realized structural substitution} from \(\Gamma\) to
\((\Delta,\sigma)\) when three conditions hold:
\begin{enumerate}
\item structural substitution transports every binding of \(\Gamma\) to its
  substituted type in \(\Delta\);
\item \(\rho\) is a homomorphism from the source structural path relation to
  \(\runtimeeq\); and
\item each source variable is mapped to a path that resolves to an endpoint
  realizing the substituted context entry.
\end{enumerate}
The theorem has mutually proved path and subtyping parts.

\begin{lemma}[Mapped substitution]
\label{lem:mapped}
Suppose \(\rho\) is a realized structural substitution from \(\Gamma\) to
\((\Delta,\sigma)\).
\begin{enumerate}
\item If \(\Gamma\struct p::\tau\), then there is an endpoint \(e\) such that
  \(\rho(p)\) resolves to \(e\) and
  \(e\in\llbracket\rho(\tau)\rrbracket_{\Delta,\sigma}\).
\item If \(\Gamma\struct\tau_1<:\tau_2\), then
  \(\rho(\tau_1)\semmap\rho(\tau_2)\).
\end{enumerate}
\end{lemma}

Most cases follow the corresponding source constructor. Selection uses the
maps stored in interval realization. Transitivity composes finite map codes.
The singleton pair case uses co-resolution and the explicit opened premise,
as described in \cref{sec:pairs}. The identity substitution and an exact
store instantiate \cref{lem:mapped} to obtain total generalized resolution
and a semantic map for every structural subtyping derivation.

\subsection{Canonical forms, preservation, and progress}

Apply a semantic map to the realization of the location singleton
\(\sing{x}\). If the target has function type, inversion of
\(\mathsf{Possible}\) produces an abstraction stored at \(x\). It also
retains the exact closure signature and the contravariant domain and
covariant codomain residuals required by beta preservation. The corresponding
pair result produces a stored pair and is used to resolve projections and
selections.

\begin{lemma}[Exact-store canonical forms]
\label{lem:canonical}
Let \(\sigma\) be exact for \(\Gamma\).
\begin{enumerate}
\item If \(\Gamma\struct\sing{x}<:(z:S)\to T\), then
  \(\sigma(x)\) is an abstraction.
\item If \(\Gamma\struct\sing{x}<:\pairty{z}{S}{a}{\tau}\), then
  \(\sigma(x)\) is a pair.
\end{enumerate}
\end{lemma}

Progress is then by induction on structural term and continuation typing.
The application case uses the function part of
\cref{lem:canonical}; path progress uses the pair part. Preservation proceeds
by cases on the six machine transitions. Beta reduction uses the retained
function-signature residuals. The let-allocation case extends the exact store
and hence the intrinsic scope; all other cases preserve the current scope.
Iteration gives \cref{thm:preservation}, and progress at its endpoint gives
\cref{thm:safety}.

\section{Mechanization and scope}
\label{sec:mechanization}

The development is checked with Lean 4 \citep{demoura2021lean4}. Syntax,
contexts, stores, and states are indexed by scope size. Renaming, weakening,
simultaneous path substitution, and opening are explicit. The default Lake
target \texttt{LambdaP} contains the present calculus and its unconditional
safety proof. The separate
\texttt{LambdaPHistory} target retains the earlier calculus, failed proof
invariants, and the checked source counterexample.

The source-facing theorems are:
\begin{center}
\scriptsize
\begin{tabular}{@{}ll@{}}
\toprule
Lean declaration & Result\\
\midrule
\texttt{Tm.Ty.closed\_progress} & initial-state progress\\
\texttt{Tm.Ty.closed\_step\_preservation} & one step from the initial state\\
\texttt{Tm.Ty.closed\_finite\_preservation} & finite-run preservation\\
\texttt{Tm.Ty.closed\_finite\_safety} & preservation and progress\\
\texttt{Tm.Ty.closed\_type\_safety} & finite-run non-stuckness\\
\bottomrule
\end{tabular}
\end{center}
They are exported from \texttt{LambdaP/Safety.lean}. The imported default
\texttt{LambdaP} target uses no \texttt{sorry}, \texttt{admit}, or
user-declared axioms.

The main proof components can be located directly:
\begin{center}
\scriptsize
\begin{tabular}{@{}ll@{}}
\toprule
Declaration & Role\\
\midrule
\texttt{Path.reduce\_progress\_precise} & source path progress\\
\texttt{Path.RuntimeEq} & co-resolution congruence\\
\texttt{Store.StructPreciseTy} & exact store invariant\\
\texttt{Path.StructCheck.mapped\_subst} & path fundamental lemma\\
\texttt{Tau.StructSub.mapped\_subst} & subtyping fundamental lemma\\
\texttt{mappedPreciseStructSafetyLaws} & canonical-form laws\\
\bottomrule
\end{tabular}
\end{center}

The theorem does not assert termination: it quantifies over endpoints of
finite executions. It does not establish algorithmic typing, type inference,
or admissibility of transitivity. Path resolution is proved deterministic;
machine-step determinism is not packaged as a theorem. Intermediate states
are preserved by the exact structural invariant, not by a literal extension
of the source term judgment.

\section{Related work}

The original DOT foundations isolate type members and path-dependent types in
an object calculus \citep{amin2014foundations}. Subsequent soundness proofs
handle substantially richer combinations of recursive objects, refinements,
and lattice structure \citep{rompf2016dot}. The tight-typing proof of
\citet{rapoport2017simple} separates reasoning in general contexts from
reasoning in restricted contexts where abstract members can be inverted.
\lambdap{} does not use inert contexts or tight typing. Its exact store and
semantic maps are tailored to immutable pair chains and the much smaller
source language.

pDOT extends DOT with paths of arbitrary length and singleton types, and
gives a Coq-mechanized safety proof \citep{rapoport2019pdot}. It is the
closest point of comparison for the distinction between path equality and
abstract selection. The present calculus differs in its dependent-pair
representation of record chains, big-step missed-label lookup, and absence of
recursive objects, intersections, and unions. The aim here is not to replace
the pDOT metatheory, but to determine which proof obligations remain when
paths and dependent members are separated from those features.

\section{Conclusion}

\lambdap{} gives a small setting for paths through dependent pairs with
abstract type members. The checked counterexample shows why singleton
symmetry cannot treat membership in an abstract selection as path identity.
This calculus makes the two roles syntactically separate. With that
distinction and pair-member subtyping restricted to singleton first
components, a direct proof establishes finite-run type safety. The open
question left by this version is precise: whether the unrestricted dependent
pair rule can be justified for the syntax-separated calculus, or whether it
admits a counterexample not captured by the present semantic invariant.

\bibliographystyle{ACM-Reference-Format}
\bibliography{references}

\end{document}
